\chapter{Informacja Fishera, efektywność i granice dokładności estymacji}
\label{ch:fisher-information}
\paragraph{Powiązania i odsyłacze.} Korzystamy z funkcji score z rozdziału~\ref{ch:likelihood} oraz z teorii estymatorów w rozdziale~\ref{ch:estimators}; przybliżenia asymptotyczne omawia rozdział~\ref{ch:asymptotic-estimation}.


\section{Motywacja: jak rozróżnialne są pobliskie parametry?}
Rozważmy dwa modele Bernoulliego: $p=0{,}50$ oraz $p=0{,}51$.
Jedna obserwacja nie pozwala ich wiarygodnie rozdzielić, nawet jeżeli
znamy dokładnie formułę rozkładu. Podobnie dwie wartości średniej
rozkładu normalnego są trudniejsze do rozróżnienia, gdy wariancja
pomiaru jest duża. \emph{Informacja Fishera} formalizuje lokalną
wrażliwość rozkładu obserwacji na zmianę parametru; nie jest miarą
informacji zawartej w pojedynczym, już zaobserwowanym wyniku.

\begin{definicja}[Score i informacja Fishera]\label{def:fisher-score}
Niech $\Theta\subset\R$ będzie przedziałem otwartym, a
$\{P_\theta:\theta\in\Theta\}$ rodziną rozkładów zdominowaną
przez miarę $\sigma$-skończoną $\mu$. Załóżmy, że istnieje wspólny zbiór $S\in\mathcal A$
taki, że $p_\theta(x)>0$ dla $x\in S$, $p_\theta(x)=0$
poza $S$ dla wszystkich $\theta$, oraz dla $\mu$-prawie
każdego $x\in S$ funkcja $\theta\mapsto p_\theta(x)$
jest różniczkowalna. Załóżmy ponadto, że
$s_\theta(x):=\partial_\theta\log p_\theta(x)$ należy do
$L^2(P_\theta)$. Funkcję $s_\theta$ nazywamy \emph{score}, a
$I(\theta):=\E_\theta[s_\theta(X)^2]$ \emph{informacją Fishera
w jednej obserwacji}. Informację w próbie zapisujemy $I_n(\theta)$.
\end{definicja}

Nie zakładamy automatycznie, że $\E_\theta s_\theta=0$.
Tę własność trzeba uzasadnić regularnością modelu. Nie należy też
mylić losowej \emph{informacji obserwowanej}
$j_\theta(X)=-\partial_\theta^2\log p_\theta(X)$ z jej
wartością oczekiwaną: dla ustalonej obserwacji wartości mogą być
bardzo różne.

\begin{twierdzenie}[Tożsamość score i addytywność informacji]
\label{thm:score-information}
Niech $p_\theta(x)$ będzie rodziną gęstości względem wspólnej
miary $\mu$ na przedziale parametrów $\Theta$. Załóżmy, że
$\partial_\theta p_\theta$ istnieje $\mu$-p.n. oraz że
$\partial_\theta\int_{\mathsf X} p_\theta(x)\,\dd\mu(x)
=\int_{\mathsf X}\partial_\theta p_\theta(x)\,\dd\mu(x)$; ponadto
$\int_{\mathsf X}(\partial_\theta p_\theta(x))^2/p_\theta(x)\,\dd\mu(x)<\infty$.
Wtedy $\E_\theta s_\theta=0$. Jeśli $X_1,\ldots,X_n$ są
niezależne o tym samym prawie $P_\theta$, to
$I_n(\theta)=nI(\theta)$.
Jeśli dodatkowo istnieje $\partial_\theta^2p_\theta$,
można dwukrotnie różniczkować całkę normalizującą i
$\E_\theta|\partial_\theta^2\log p_\theta(X)|<\infty$,
to $I(\theta)=-\E_\theta\partial_\theta^2\log p_\theta(X)$.
\end{twierdzenie}
\begin{proof}
Z definicji $s_\theta p_\theta=\partial_\theta p_\theta$,
więc $\E_\theta s_\theta=\int_{\mathsf X}\partial_\theta p_\theta(x)
\,\dd\mu(x)=\partial_\theta1=0$. Dla próby score jest sumą
$S_n=\sum_i s_\theta(X_i)$: wynika to z logarytmowania
iloczynu gęstości. Po podniesieniu do kwadratu
$\E S_n^2=\sum_i\E s_\theta(X_i)^2+
2\sum_{i<j}\E s_\theta(X_i)s_\theta(X_j)$.
Niezależność i zerowanie średniej eliminują wszystkie wyrazy
mieszane, zatem $I_n=nI$.
Z równości $\partial_\theta^2\log p=(\partial_\theta^2p)/p-
(\partial_\theta p)^2/p^2$ dostajemy
$\E_\theta\partial_\theta^2\log p=
\int_{\mathsf X}\partial_\theta^2p_\theta(x)\,\dd\mu(x)-I(\theta)=-I(\theta)$.
Ostatnia równość wykorzystuje dopuszczoną zamianę drugiej
pochodnej z całką i nie jest poprawna bez niej.
\end{proof}

\begin{zadanie}[Bernoulli: dlaczego skrajne prawdopodobieństwa są szczególne?]
Niech $X\sim\operatorname{Bernoulli}(p)$, $p\in(0,1)$.
Oblicz score i informację, a następnie podaj przybliżony
błąd standardowy $\widehat p$ z $n$ niezależnych prób.
\par\emph{Rozwiązanie.} Mamy
$\ell(p;x)=x\log p+(1-x)\log(1-p)$,
$s_p(x)=x/p-(1-x)/(1-p)=(x-p)/(p(1-p))$.
Ponieważ $\Var_p X=p(1-p)$, informacja wynosi
$I(p)=1/[p(1-p)]$ i $I_n(p)=n/[p(1-p)]$.
Dla $\widehat p=\bar X$ zachodzi dokładnie
$\Var_p\widehat p=p(1-p)/n$. Stąd po podstawieniu
$\widehat p$ orientacyjny błąd standardowy to
$\sqrt{\widehat p(1-\widehat p)/n}$.
Blisko $p=0$ lub $1$ wzór wymaga ostrożności: punkty brzegowe
nie należą do otwartej przestrzeni parametrów, a dla $\widehat p=0$
nie można wnioskować, że niepewność parametru faktycznie znika.
\end{zadanie}

\begin{przyklad}[Poisson z ekspozycją i normalny pomiar]
Jeśli niezależne $Y_i\sim\operatorname{Pois}(\lambda t_i)$,
$t_i>0$, to score jest $S(\lambda)=\sum_iY_i/\lambda-
\sum_it_i$, a informacja w całej próbie wynosi
$I(\lambda)=\sum_it_i/\lambda$. To \emph{łączny czas ekspozycji},
a nie liczba okresów, wyznacza podstawową skalę dokładności
estymacji intensywności. Jeśli $X_i\sim\Normal(\mu,\sigma^2)$
przy znanym $\sigma$, to $s_\mu(X_i)=(X_i-\mu)/\sigma^2$,
$I_n(\mu)=n/\sigma^2$: większy szum oznacza mniejszą informację.
\end{przyklad}

\section{Nierówność Craméra--Rao i warunki jej stosowania}
\begin{twierdzenie}[Nierówność Craméra--Rao dla funkcji parametru]
\label{thm:cramer-rao}
Niech $\Theta\subset\R$ będzie przedziałem otwartym,
$T$ rzeczywistą statystyką z $\E_\theta T^2<\infty$,
a $g(\theta)=\E_\theta T$. Załóżmy, że score
$S_\theta=\partial_\theta\log p_\theta(X)$ całej obserwacji
jest kwadratowo całkowalny, $\E_\theta S_\theta=0$,
$0<I_n(\theta)<\infty$ oraz że wolno różniczkować
$\int_{\mathsf X} T(x)p_\theta(x)\,\dd\mu(x)$ pod całką. Wtedy
\[
\Var_\theta(T)\ge\frac{[g'(\theta)]^2}{I_n(\theta)}.
\]
Równość dla ustalonego $\theta$ zachodzi (gdy $g'(\theta)\ne0$)
dokładnie wtedy, gdy $T-g(\theta)$ jest $P_\theta$-p.n.
proporcjonalna do $S_\theta$.
\end{twierdzenie}
\begin{proof}
Z dozwolonego różniczkowania i tożsamości
$\partial_\theta p_\theta=S_\theta p_\theta$ otrzymujemy
$g'(\theta)=\E_\theta[T S_\theta]$.
Ponieważ $\E_\theta S_\theta=0$, ta wartość jest kowariancją
$\Cov_\theta(T,S_\theta)$. Nierówność Cauchy'ego--Schwarza
daje $|g'(\theta)|^2\le\Var_\theta(T)\Var_\theta(S_\theta)$.
Ostatnia wariancja równa się $I_n(\theta)$, co po podzieleniu
przez dodatnią liczbę daje tezę. Warunek równości w
Cauchym--Schwarzu jest równoważny liniowej zależności
zmiennych scentrowanych w $L^2(P_\theta)$.
\end{proof}

\begin{wniosek}[Znak równości w modelu normalnym]
Jeśli $X_1,\ldots,X_n\iid\Normal(\mu,\sigma^2)$ przy
znanej $\sigma^2>0$, to dla nieobciążonego estymatora
$\bar X$ średniej $\mu$ mamy
$\Var_\mu(\bar X)=\sigma^2/n=1/I_n(\mu)$.
Zatem $\bar X$ osiąga granicę Craméra--Rao przy każdym $\mu$.
\end{wniosek}

\begin{uwaga}[Czemu rozkład jednostajny łamie naiwny rachunek?]
Dla próby $X_i\iid U(0,\theta)$ nośnik $(0,\theta)^n$
zależy od $\theta$. Na wnętrzu nośnika log-wiarygodność ma
pochodną $-n/\theta$, której wartość oczekiwana jest
$-n/\theta\ne0$. Zatem krok $\E_\theta S_\theta=0$
nie zachodzi i poprzednie twierdzenie nie ma zastosowania.
Co więcej, skorygowany estymator maksimum ma wariancję rzędu
$n^{-2}$, a nie $n^{-1}$. Nie jest to paradoks, lecz efekt
informacji pochodzącej z \emph{brzegu nośnika}.
\end{uwaga}

\section{Wiele parametrów i parametry zakłócające}
\begin{twierdzenie}[Macierzowa nierówność Craméra--Rao]
\label{thm:cr-matrix}
Niech $\Theta\subset\R^d$ będzie otwarta, a score całej próby
$S_\theta=\nabla_\theta\log p_\theta(X)$ spełnia
$\E_\theta S_\theta=0$, $I_n(\theta)=\E_\theta[S_\theta
S_\theta^{\mathsf T}]$ jest dodatnio określona oraz
$\nabla_\theta\E_\theta T=\E_\theta[T S_\theta]$ dla
$T\in L^2(P_\theta)$. Oznaczmy $g(\theta)=\E_\theta T$.
Wtedy
\[
\Var_\theta(T)\ge \nabla g(\theta)^{\mathsf T}
I_n(\theta)^{-1}\nabla g(\theta).
\]
\end{twierdzenie}
\begin{proof}
Niech $a=I_n^{-1}\nabla g$ i $Z=T-g-a^{\mathsf T}S_\theta$.
Ponieważ $\E Z=0$ oraz $\E[(T-g)S_\theta]=\nabla g$,
rozwinięcie kwadratu i $\E[S_\theta S_\theta^{\mathsf T}]=I_n$
dają $\E Z^2=\Var(T)-2a^{\mathsf T}\nabla g+a^{\mathsf T}I_na
=\Var(T)-\nabla g^{\mathsf T}I_n^{-1}\nabla g$.
Nieujemność $\E Z^2$ kończy dowód.
\end{proof}

Jeśli $\theta=(\psi,\eta)$, gdzie $\psi$ jest interesującym
nas parametrem, a $\eta$ parametrem zakłócającym, należy użyć
pełnej macierzy informacji. Dla odpowiedniego rozkładu blokowego
$I=\begin{psmallmatrix}I_{\psi\psi}&I_{\psi\eta}\\
I_{\eta\psi}&I_{\eta\eta}\end{psmallmatrix}$, gdy dolny blok
jest odwracalny, pierwszy element $I^{-1}$ to
$(I_{\psi\psi}-I_{\psi\eta}I_{\eta\eta}^{-1}I_{\eta\psi})^{-1}$.
Wyrażenie w nawiasie nazywamy informacją efektywną o $\psi$.
Nie wolno więc automatycznie używać $1/I_{\psi\psi}$,
jak gdyby $\eta$ było znane.

\begin{zadanie}[Normalny model z nieznaną wariancją]
Dla próby $X_i\iid\Normal(\mu,v)$, $\mu\in\R$, $v>0$,
oblicz macierz informacji parametrów $(\mu,v)$.
\par\emph{Rozwiązanie.} Pojedynczy score ma składowe
$s_\mu=(X-\mu)/v$ i
$s_v=-1/(2v)+(X-\mu)^2/(2v^2)$.
Z symetrii rozkładu normalnego składnik mieszany wynosi zero.
Z momentów centralnych $\E(X-\mu)^2=v$ i
$\E(X-\mu)^4=3v^2$ otrzymujemy
$I_1=\diag(1/v,1/(2v^2))$.
Dla $n$ niezależnych pomiarów $I_n=nI_1$.
Zatem parametr $v$ jest tu informacyjnie ortogonalny do $\mu$:
nieznana wariancja nie zmienia elementu $(I_n^{-1})_{\mu\mu}$,
choć istotnie zmienia sposób szacowania błędu standardowego z danych.
\end{zadanie}


\section{Lokalna geometria modeli: informacja a dywergencja KL}
Tożsamość informacyjna ma dodatkową interpretację, która jest
przydatna przy projektowaniu doświadczenia. Dwa bardzo bliskie
parametry nie dają dwóch zupełnie różnych rozkładów; w modelach
regularnych ich odległość w sensie dywergencji KL jest
\emph{drugiego rzędu} względem różnicy parametrów. Informacja
Fishera jest współczynnikiem tego rozwinięcia.

\begin{stwierdzenie}[Lokalne rozwinięcie dywergencji KL]
\label{prop:kl-fisher-expansion}
Niech $\Theta\subset\R$ będzie otwartym przedziałem, a
$\{p_\eta:\eta\in\Theta\}$ modelem zdominowanym o dodatnich
gęstościach na wspólnym nośniku. Ustalmy $\theta\in\Theta$.
Załóżmy, że funkcja
$K(h)=\E_\theta[\log p_\theta(X)-\log p_{\theta+h}(X)]$
jest dwukrotnie różniczkowalna w otoczeniu $h=0$, przy czym
wolno różniczkować pod wartością oczekiwaną, a tożsamości
informacyjne obowiązują dla $\eta$ w tym otoczeniu.
Wtedy
$K(h)=\tfrac12 I(\theta)h^2+o(h^2)$, gdy $h\to0$.
\end{stwierdzenie}
\begin{proof}
Ponieważ $K(0)=0$ i $K'(0)=-\E_\theta s_\theta(X)=0$,
wystarczy policzyć drugą pochodną w zerze.
Mamy $K''(0)=-\E_\theta\partial_\theta^2\log p_\theta(X)
=I(\theta)$. Rozwinięcie Taylora z resztą Peana daje tezę.
Przytoczone założenia są istotne: przy nośniku zależnym od
parametru nawet wartość $K(h)$ może być nieskończona w jednej
stronie zera.
\end{proof}

\begin{przyklad}[Dlaczego powtarzanie pomiaru pomaga?]
Dla $P_\mu=\Normal(\mu,\sigma^2)$ przy znanym $\sigma>0$
rachunek gaussowski daje dokładnie
$D_{\rm KL}(P_\mu\|P_{\mu+h})=h^2/(2\sigma^2)$.
Dla $n$ niezależnych pomiarów dywergencja praw wspólnych jest
$n$ razy większa. Dlatego przy ustalonym stosunku sygnału
do szumu wzrost liczby niezależnych pomiarów zwiększa
rozróżnialność parametrów. Jeśli jednak wielokrotnie mierzymy
ten sam obiekt przy wspólnym błędzie systematycznym, model
iloczynowy nie jest poprawny i wniosek o addytywności nie musi
pozostać prawdziwy.
\end{przyklad}

\section{Efektywność i ryzyko: co rzeczywiście mierzymy?}
W statystyce mówi się czasem, że estymator jest \emph{efektywny},
gdy osiąga granicę Craméra--Rao w danym regularnym modelu.
To jest własność \emph{względem konkretnej funkcji straty i klasy
estymatorów}, a nie stwierdzenie, że jest to najlepsza procedura
do wszystkich zastosowań. Granica Craméra--Rao dotyczy w podanej
wersji estymatorów o ustalonej wartości oczekiwanej $g(\theta)$.
Estymator obciążony może mieć mniejsze MSE; metody odporne mogą
być lepsze, gdy część pomiarów jest zanieczyszczona.

\begin{zadanie}[Dwa sposoby estymacji wariancji normalnej]
Dla $X_i\iid\Normal(\mu,\sigma^2)$ o nieznanych $\mu$ i
$\sigma^2$ porównaj estymator nieobciążony
$S^2=(n-1)^{-1}\sum_i(X_i-\bar X)^2$ i MLE
$V_n=n^{-1}\sum_i(X_i-\bar X)^2$ przy stracie kwadratowej.
\par\emph{Rozwiązanie.} Z rozkładu
$(n-1)S^2/\sigma^2\sim\chi^2_{n-1}$ wynika
$\E S^2=\sigma^2$ i $\Var S^2=2\sigma^4/(n-1)$.
Dla $V_n=(n-1)S^2/n$ mamy
$\E V_n=(n-1)\sigma^2/n$ i
$\Var V_n=2(n-1)\sigma^4/n^2$. Dlatego
\[
\MSE_\theta(V_n)=\frac{2(n-1)+1}{n^2}\sigma^4
=\frac{2n-1}{n^2}\sigma^4,
\qquad
\MSE_\theta(S^2)=\frac{2}{n-1}\sigma^4.
\]
Dla $n\ge2$ pierwszy z tych błędów jest mniejszy.
Ten przykład przypomina, że minimalna wariancja w klasie
estymatorów nieobciążonych nie rozstrzyga optymalności przy
uwzględnieniu estymatorów obciążonych.
\end{zadanie}

\section{Kiedy granica informacyjna daje praktyczną odpowiedź?}
\begin{center}\small
\begin{tabular}{>{\raggedright\arraybackslash}p{0.25\textwidth}>{\raggedright\arraybackslash}p{0.32\textwidth}>{\raggedright\arraybackslash}p{0.30\textwidth}}
\toprule
\textbf{Cel} & \textbf{Co można obliczyć?} & \textbf{Czego nie wolno wnioskować?}\\
\midrule
Porównanie dwóch nieobciążonych estymatorów & Wariancje, MSE i dolną granicę Craméra--Rao & Osiąganie granicy nie jest gwarantowane\\
Projektowanie pomiaru & Informację na obserwację lub ekspozycję & Więcej rekordów nie zawsze znaczy więcej informacji przy zależności\\
Parametr z nuisance & Macierz informacji i dopełnienie Schura & Nie wolno pominąć współparametrów\\
Mała próba i punkt brzegowy & Dokładny rozkład estymatora, przedział dokładny & Asymptotyczna normalność może być złym przybliżeniem\\
Model o zmiennym nośniku & Bezpośredni rachunek rozkładu statystyki & Nie stosować regularnego Craméra--Rao bez założeń\\
\bottomrule
\end{tabular}
\end{center}

\begin{zadanie}[Jednostki informacji przy obserwacji zdarzeń]
Dwie maszyny pracowały odpowiednio $20$ i $80$ godzin;
zaobserwowano łącznie $7$ awarii. W modelu niezależnych zliczeń
Poissona o wspólnej intensywności $\lambda$ oblicz MLE i
przybliżony błąd standardowy.
\par\emph{Rozwiązanie.} Łączna ekspozycja to $T=100$ godzin,
więc $\widehat\lambda=7/100=0{,}07$ awarii na godzinę.
Obserwowana krzywizna log-wiarygodności przy maksimum wynosi
$j(\widehat\lambda)=7/\widehat\lambda^2$;
informacja Fishera to $T/\lambda$.
Oba sposoby po podstawieniu MLE prowadzą do
$\widehat{\mathrm{se}}=\sqrt{\widehat\lambda/T}
\approx0{,}0265$ na godzinę. Jest to oszacowanie asymptotyczne:
przy zaledwie siedmiu zdarzeniach warto porównać je z dokładną
procedurą opartą na rozkładzie Poissona.
\end{zadanie}

% ================================================================
% ROZDZIAL 8 -- ASYMPTOTYCZNA TEORIA ESTYMACJI
% ================================================================
